% =============================================================================
\chapter{Métricas de Avaliação}
% =============================================================================

\section{Desempenho da tarefa}

\begin{itemize}
\item taxa de sucesso da missão e sobrevivência;
\item conclusão de objetivos e sobrevivência de aliados;
\item dano recebido e recursos consumidos;
\item duração do episódio e número de replanejamentos.
\end{itemize}

\section{Desempenho multiobjetivo}

As métricas serão selecionadas conforme as hipóteses do cenário e reportadas
com suas referências, em particular o protocolo de benchmarking MORL
\cite{felten2023toolkit}.

\begin{itemize}
\item retorno vetorial acumulado;
\item utilidade escalarizada sob diferentes preferências;
\item hipervolume e cobertura de soluções somente quando houver conjunto de soluções
e ponto de referência explicitamente definidos;
\item utilidade esperada sobre uma distribuição de preferências;
\item \textit{Maximum Utility Loss} (MUL), quando houver solução de referência;
\item \textit{Inverted Generational Distance} (IGD), apenas em domínios com
fronteira de Pareto de referência conhecida;
\item desempenho sob vetores de preferência não vistos.
\end{itemize}

Para utilidade linear, a utilidade esperada será reportada apenas quando a
distribuição sobre preferências estiver explicitamente definida; métricas de
hipervolume não substituem a avaliação sob preferências condicionadas.

\section{Adaptação e generalização}

\begin{itemize}
\item desempenho após mudanças contextuais;
\item latência de adaptação após alteração de preferências;
\item frequência e adequação das mudanças de método;
\item recuperação após eventos inesperados;
\item desempenho em cenários e combinações não observados no treinamento.
\end{itemize}

\section{Diversidade e coerência comportamental}

\begin{itemize}
\item distribuição dos métodos HTN selecionados;
\item número de estratégias distintas observadas;
\item variação entre personalidades e vetores de preferência;
\item consistência sob estados e preferências equivalentes;
\item taxa de oscilação ou troca prematura de estratégia.
\end{itemize}

\section{Validade semântica}

\begin{itemize}
\item taxa de seleção de método inválido;
\item violações de pré-condições;
\item tentativas de uso de recursos indisponíveis;
\item violações de sequência lógica de tarefas.
\end{itemize}

Para a arquitetura completa, espera-se taxa de seleção inválida igual a zero,
pois os métodos não aplicáveis são removidos antes da decisão MORL. Falhas de
decomposição, falhas de execução e replanejamentos devem ser contabilizados
separadamente.

\section{Desempenho computacional}

\begin{itemize}
\item latência total de decisão, tempo da fonte de preferências e tempo de inferência;
\item tempo de filtragem HTN;
\item tempo de decomposição do método selecionado;
\item taxa e custo de \emph{method fallback};
\item tempo de treinamento;
\item quantidade de métodos decompostos no caminho normal;
\item consumo de memória;
\item estabilidade da frequência de decisão em tempo real.
\end{itemize}

\section{Densidade de decisão}

Além de medir latência, a avaliação reportará a densidade de decisão: número de
consultas ao seletor por episódio, por segundo simulado e por ação primitiva.
Essa métrica testa se a abstração HTN reduz a frequência efetiva de escolhas sem
presumir, a priori, que seu espaço de decisão é globalmente menor.

\section{Rastreabilidade e auditabilidade}

Cada decisão será registrada com tarefa composta, métodos aplicáveis, métodos
rejeitados e pré-condições falhas, vetor de preferências, valores objetivos,
utilidade escalarizada e método escolhido.
A avaliação verificará completude dos traços e coerência entre prioridades e
decisão, complementada por inspeção qualitativa de episódios representativos.
